\section{Benchmark Construction}

\subsection{Pipeline and data layers}
Supplementary Figure~\ref{fig:overview} summarizes the artifact. Each of 50,000 utility-IID clusters contains four source-role variants, producing 200,000 rows split by cluster into 120,000 training, 40,000 validation, and 40,000 test rows. Before outcome inspection, the 10,000 test clusters were ordered by a SHA-256 rule and split exactly into 5,000 paper-test and 5,000 community-hidden clusters. The seven frozen opportunity slices include ordinary, moderate stress, severe stress, rare safe opportunity, high missed opportunity, stress-only period, and a no-opportunity control.

\paragraph{Branch-neutral construction.} The generator follows a fixed sequence:
\begin{enumerate}
  \item sample an event state, opportunity intent, stress severity, liquidity, and cost regime;
  \item sample four source-role-specific noisy priors and confidence/uncertainty fields;
  \item realize return, full/reduced post-cost utility, and outcome-only harm labels;
  \item attach provenance transformations or task-specific workflow state; and
  \item run execution rules only after the row and legal-field manifest are frozen.
\end{enumerate}
An oracle-positive opportunity indicates that at least one role- and provenance-eligible full or reduced action has positive post-cost utility. It is never a deployable feature. Generator parameters, expected opportunity rates, mechanism holdouts, and seeds are frozen before evaluation; validation defects remain in the negative-results log and were repaired only by rule-independent corrections.

To test generator dependence, we prospectively add two mechanisms without changing rule thresholds: a sequential-market generator with regimes, inventory, market impact, and path dependence, and an institutional-workflow generator with trading, credit, and research-to-action tasks, risk budgets, compliance flags, approval depth, and review pressure. Together they add 96,000 rows from 24,000 clusters. They are mechanistically distinct but still author-controlled scenarios, not market simulators or institutional logs.

\subsection{Provenance layer}
The provenance layer reuses the controlled market-state scaffold but adds role claims, verified roles, original source identifiers, delegation parents, transformation chains, content hashes, hop depth, traceability, and lineage attestation. Five attack families are balanced at the cluster level: clean, role noise, delegation, paraphrase, and multi-hop. Role noise rates are 0\%, 5\%, 10\%, 20\%, and 40\%; hop depths range from one to five. Traceable rows expose enough lineage for detection in principle. Untraceable rows intentionally omit that evidence and are reported as a separate stress/impossibility stratum.

\subsection{Point-in-time public, order-book, and training layers}
The public layer contains 2,500 structurally valid Polymarket event clusters, split chronologically into 1,500 training, 500 validation, and 500 sealed test clusters. The source probability is the last observation strictly before a fixed decision timestamp at 35\% of scheduled lifecycle. Outcome and post-decision fields are labels only. Twenty-three high-activity months reached the frozen pagination cap and are disclosed as truncated. The pre-result power gate estimates 0.164 power, with LCB95 0.160, against a 0.80 requirement, so the layer remains \texttt{exploratory\_only}. FRED and GDELT remain documented extensions rather than confirmatory sources.

A v0.3 extension adds observed market-outcome replay with constructed priors and roles: 185,592 Binance rows from 703 UTC dates and 16,776 licensed MES rows from 233 RTH sessions. Binance combines official banded depth with one-minute BTCUSDT/ETHUSDT bars, split 120/500/83 into development, confirmatory paper test, and hidden dates; planning power uses the frozen v0.2 variance model. MES uses one-contract BBO replay split 60/140/33, with the paper test descriptive only. Both sources reuse frozen priors, rules, and thresholds; outcomes remain sealed until one-time evaluation, and MES rows are not redistributed.

A v0.6 actual-proposal extension uses 300 non-overlapping dates from May 2025 through March 2026, deterministically assigned across BTC, ETH, BNB, SOL, and XRP. Each date contains a directional-execution task with an eligible edge proposer and a 25\% risk-limit-increase task whose source role is deterministically assigned as eligible proposer or evidence-only critic before generation. Three model configurations receive the same legal pre-decision context. The chronological split contains 50 development, 200 one-time paper-test, and 50 author-unevaluated holdout dates. Holdout proposals were generated and encrypted before outcome evaluation; their key remains outside the repository and the ciphertext was not decrypted. Outputs retain malformed placeholders without repair or favorable reruns. Rationales are audit metadata and are excluded from every rule.

The registered training layer freezes D0--D7 at 400 clusters per seed and five seeds. A primary endpoint is valid only when every held-out mechanism has a 95\% coverage lower bound of at least 0.05. A prospective v0.3 robustness study preserves that primary and adds Logistic Regression, histogram gradient boosting, MLP, and selective Logistic learners; 400, 1,000, 5,000, and 20,000-cluster budgets; three curricula; and six held-out mechanisms. The validity condition is deliberately allowed to make a comparison undefined.

\begin{table*}[t]
  \caption{\textbf{FinAuth-Audit releases multiple validity layers rather than a single aggregate row count.} Counts and boundaries are generated from manifests. Controlled, provenance, external order-book, and actual-model results use separately frozen one-time paper tests; the event-market and mechanistic/training extensions retain their stated validation-only or sealed status.}
  \label{tab:dataset_summary}
  \centering
  \benchmarktableformat
  % Generated by finauth_audit/paper/generate_tables.py. Do not edit manually.
\begin{tabular}{p{0.17\linewidth}>{\raggedleft\arraybackslash}p{0.10\linewidth}>{\raggedleft\arraybackslash}p{0.10\linewidth}p{0.23\linewidth}p{0.28\linewidth}}
\toprule
Layer & Rows & Clusters & Split or budget & Coverage / boundary \\
\midrule
Controlled core & 200,000 & 50,000 & 30k / 10k / 5k / 5k clusters & train / val. / paper / hidden \\
Provenance laundering & 200,000 & 50,000 & 30k / 10k / 5k / 5k clusters & 5 attacks; 2 trace strata \\
Mechanistic generators & 96,000 & 24,000 & sequential / institutional & validation robustness; test sealed \\
Public point-in-time & 2,500 & 2,500 & 1.5k / 0.5k / 0.5k & 500 validation reported; test sealed \\
External order book & 202,368 & 936 & Binance 120/500/83; MES 60/140/33 & powered public + licensed descriptive \\
Actual model proposals & 1,200 & 200 & 50 / 200 / 50 dates & 3 models; 2 tasks/date; aggregate only \\
Training utility & 144 fits & 400--20k / curriculum & 3 curricula; 4 learners & 6 mechanisms; N/A retained \\
\bottomrule
\end{tabular}

\end{table*}
